We have demonstrated that the bisect recursive bisection algorithm extends
naturally to multi-member district generation by adjusting the district-count
parameter $k$ and adapting the per-seat balance tolerance formula.  Across
five configurations in three states, MMDs systematically improve minority
representation relative to enacted single-member plans---California's
4-member OLPR districts produce an estimated 22 minority-opportunity seats
versus 17 in the enacted plan---while incurring a modest compactness penalty
of approximately 5--12\% in mean Polsby-Popper depending on district
magnitude.\singrun

\paragraph{STV vs.\ OLPR.}
The two allocation rules produce similar seat distributions.  OLPR produces
one additional minority-opportunity seat in the California case study, driven
by the direct candidate-preference mechanism.  STV has the advantage of
requiring no party registration, making it more accessible for jurisdictions
with strong traditions of cross-party voting.

\paragraph{Proportionality-compactness tradeoff.}
The tradeoff is real but modest: the 4-member configuration achieves
Gallagher indices of 1.8--2.9 (comparable to Nordic PR systems) at a 10--12\%
compactness cost.  The 2-member configuration achieves intermediate results
in both dimensions.  The choice between magnitudes is primarily a values
judgment---how much compactness to sacrifice for proportionality---rather
than an algorithmic question.

\paragraph{Legal path.}
The statutory barrier to MMDs (2 U.S.C.\ \S~2c) is a legislative choice, not
a constitutional requirement.  The Fair Representation Act would clear this
barrier and explicitly authorize the type of MMD plan generated in this paper.
The 110-year Illinois precedent with three-member Cumulative Voting districts
demonstrates that MMDs are compatible with American political culture.

\paragraph{Future work.}
Several extensions are warranted.  First, ensemble analysis (multiple seeds)
would characterise the distribution of possible MMD configurations rather than
the single-run results reported here.  Second, the full 50-state analysis
would establish national baselines for proportionality and minority
representation under various MMD designs.  Third, combined MMD and VRA
edge-weighting experiments would explore whether geographic minority clustering
and proportional allocation are complementary or substitutable strategies.
